% 
%            ,,,                                        
%          `7MM            _.o9                                
%            MM                                             
%  ,6"Yb.    MM  ,p6"bo   ,6"Yb.  M"""MMV  ,6"Yb.  `7Mb,od8 
% 8)   MM    MM 6M'  OO  8)   MM  '  AMV  8)   MM    MM' "' 
%  ,pm9MM    MM 8M        ,pm9MM    AMV    ,pm9MM    MM     
% 8M   MM    MM YM.    , 8M   MM   AMV  , 8M   MM    MM     
% `Moo9^Yo..JMML.YMbmd'  `Moo9^Yo.AMMmmmM `Moo9^Yo..JMML.   
%
%
% 神经网络教材 - 摘要

\newpage
\thispagestyle{empty}

\begin{center}
    \Large\textbf{摘要}
\end{center}

\vspace{1cm}

神经网络彻底改变了人工智能领域，使计算机视觉、自然语言处理、语音识别等众多领域取得了突破性进展。本教材全面介绍神经网络的基础知识，从基本概念到最先进的网络架构。

本书共十一章，内容包括：

\begin{enumerate}
    \item \textbf{数学基础：} 涵盖理解神经网络所需的线性代数、微积分和概率论基础知识。
    
    \item \textbf{激活函数：} 对Sigmoid、Tanh、ReLU及各种高级激活函数进行全面分析，并提供选择指南。
    
    \item \textbf{损失函数：} 详细介绍交叉熵、均方误差及各类任务专用损失函数。
    
    \item \textbf{神经网络架构：} 深度神经网络(DNN)、卷积神经网络(CNN)和循环神经网络(RNN)。
    
    \item \textbf{核心操作：} 卷积、池化和归一化操作，配有严格的数学推导。
    
    \item \textbf{高级机制：} 跳跃连接和通道注意力机制，包括残差学习。
    
    \item \textbf{优化：} 梯度下降算法、各类优化器（SGD、Adam、AdamW）和学习率调度器。
    
    \item \textbf{经典架构：} 从LeNet和AlexNet到ResNet、VGG及Transformer架构。
    
    \item \textbf{计算分析：} 网络效率分析的参数量统计和FLOPs计算。
    
    \item \textbf{训练与框架：} 实际训练过程、超参数调优、PyTorch实现及CUDA编程概念。
    
    \item \textbf{数据集处理：} 数据收集、标注、增强和划分策略。
\end{enumerate}

每章都包含详细的数学推导、实际案例以及对经典和最新研究论文的参考文献。完成本教材学习后，读者将具备设计、实现和训练各类神经网络应用的能力。

\vspace{2cm}

\textbf{关键词：} 神经网络、深度学习、机器学习、CNN、RNN、Transformer、优化、反向传播、PyTorch、计算机视觉、自然语言处理

\newpage
